\documentclass[11pt,a4paper]{article}
\usepackage{cmap}
\usepackage[T2A]{fontenc}
\usepackage[utf8]{inputenc}
\usepackage[russian]{babel}
\usepackage{amsmath,amssymb,amsthm}
\usepackage{geometry}
\geometry{left=20mm,right=20mm,top=20mm,bottom=22mm}
\usepackage{tikz}
\usetikzlibrary{arrows.meta,positioning,shapes.geometric}
\usepackage{xcolor}
\usepackage{tabularx,array,booktabs,multicol}
\usepackage{fancyhdr}
\usepackage{enumitem}
\usepackage{microtype}
\usepackage{needspace}
\usepackage{setspace}
\usepackage{hyperref}
\hypersetup{hidelinks,pdftitle={НОД и НОК},pdfauthor={Учебно-методическое пособие}}
\definecolor{ink}{HTML}{253644}
\definecolor{navy}{HTML}{324F69}
\definecolor{slate}{HTML}{798E9F}
\definecolor{muted}{HTML}{586B78}
\definecolor{pale}{HTML}{EDF2F5}
\color{ink}
\setstretch{1.075}
\setlength{\parindent}{0pt}
\setlength{\parskip}{3.8pt}
\setlist[itemize]{leftmargin=6mm,itemsep=1.5pt,topsep=3pt}
\setcounter{secnumdepth}{0}
\pagestyle{fancy}\fancyhf{}
\fancyhead[L]{\footnotesize\color{muted} ЕГЭ по математике\,|\,Целые числа}
\fancyhead[R]{\footnotesize\color{muted} НОД и НОК}
\fancyfoot[C]{\small\color{muted}\thepage}
\renewcommand{\headrulewidth}{0.35pt}
\setlength{\headheight}{14pt}
\renewcommand{\arraystretch}{1.32}
\newcommand{\gcdru}[2]{\mathop{\text{НОД}}\nolimits(#1,#2)}
\newcommand{\lcmru}[2]{\mathop{\text{НОК}}\nolimits(#1,#2)}
\newcommand{\blocktitle}[2]{\clearpage\section{Блок #1. #2}\vspace{-2mm}\noindent\color{slate}\rule{\linewidth}{.6pt}\color{ink}\vspace{2mm}}
\newcommand{\concept}[1]{\par\needspace{4\baselineskip}{\color{navy}\bfseries #1.}\enspace}
\newcommand{\example}[1]{\par\needspace{5\baselineskip}{\color{navy}\bfseries Разобранный пример.} #1\par}
\newcommand{\warning}[1]{\par\needspace{4\baselineskip}{\color{muted}\bfseries Типичная ошибка.} #1\par}
\newcommand{\selfcheck}[1]{\par\needspace{3\baselineskip}{\color{muted}\bfseries Проверка себя.} #1\par}
\newcommand{\minitrain}{\par\needspace{8\baselineskip}\vspace{1mm}{\color{navy}\bfseries Три мини-упражнения}\par\vspace{-1mm}}
\newcommand{\mini}[2]{\par\noindent\hspace*{2mm}\textbf{#1)}\enspace #2\par}
\newcommand{\task}[2]{\par\needspace{2\baselineskip}\noindent\makebox[8mm][l]{\color{navy}\bfseries #1.}#2\par\vspace{2mm}}
\newcommand{\training}[2]{\clearpage\section*{Тренировочный блок #1}\textcolor{muted}{#2}\par\vspace{3mm}\rule{\linewidth}{.4pt}\vspace{2mm}}
\newcommand{\answers}[2]{\clearpage\section*{Ответы к тренировочному блоку #1}\vspace{3mm}#2}
\newcommand{\anstable}[1]{\begin{center}\small\begin{tabularx}{\linewidth}{>{\centering\arraybackslash}p{9mm}X>{\centering\arraybackslash}p{9mm}X}\toprule\textbf{№}&\textbf{Ответ}&\textbf{№}&\textbf{Ответ}\\\midrule #1\\\bottomrule\end{tabularx}\end{center}}
\tikzset{st/.style={ellipse,draw=navy,line width=.8pt,fill=pale,align=center,minimum height=10mm,text width=45mm,font=\small},
pr/.style={rectangle,rounded corners=2mm,draw=slate,line width=.8pt,align=center,minimum height=12mm,text width=67mm,font=\small},
de/.style={diamond,aspect=2.35,draw=navy,line width=.8pt,align=center,text width=32mm,inner sep=1.8mm,font=\small},
ar/.style={-{Stealth[length=2.4mm,width=1.7mm]},line width=.8pt,draw=navy},
ann/.style={font=\footnotesize\bfseries,inner sep=1.5pt,color=navy}}
\newcommand{\schemetitle}[1]{\clearpage\section*{Алгоритм: #1}\vspace{10mm}}
\begin{document}
\begin{titlepage}
\thispagestyle{empty}
\vspace*{26mm}
{\color{navy}\large УЧЕБНО-МЕТОДИЧЕСКОЕ ПОСОБИЕ}\par
\vspace{14mm}
{\Huge\bfseries НОД и НОК}\par
\vspace{9mm}
{\Large Наибольший общий делитель\\[2mm] и наименьшее общее кратное}\par
\vspace{18mm}
{\color{muted}\large Определения, методы вычисления,\\[2mm] взаимная простота и задачи ЕГЭ}\par
\vspace{22mm}
\begin{tikzpicture}[x=1cm,y=1cm]
\draw[navy,line width=.95pt] (0,0) .. controls (3,.8) and (5.8,-.7) .. (8.3,.1) .. controls (10.5,.8) and (13,.15) .. (15.6,.5);
\draw[slate,line width=.35pt] (0,-.35) -- (15.6,-.35);
\end{tikzpicture}
\vfill
\end{titlepage}
\clearpage
\section*{Карта темы}
\textbf{Главная идея.} НОД отвечает на вопрос «какое наибольшее число делит каждое из данных чисел?», а НОК --- «какое наименьшее положительное число делится на каждое из данных чисел?». Эти понятия связаны с простыми множителями, остатками и взаимной простотой.
\begin{center}
\begin{tikzpicture}[node distance=8mm]
\node[pr,text width=116mm] (a) {\textbf{1. Смысл и вычисление}\quad общие делители и кратные, простые множители};
\node[pr,text width=116mm,below=of a] (b) {\textbf{2. Алгоритм Евклида}\quad последовательные остатки, линейные комбинации};
\node[pr,text width=116mm,below=of b] (c) {\textbf{3. Взаимная простота}\quad свойства, связь НОД и НОК, обратные задачи};
\node[pr,text width=116mm,below=of c] (d) {\textbf{4. Применение и доказательство}\quad группировка, периодичность, задачи с параметром};
\draw[ar](a)--(b);\draw[ar](b)--(c);\draw[ar](c)--(d);
\end{tikzpicture}
\end{center}
\concept{Обозначения} Здесь $\mathbb N=\{1,2,3,\ldots\}$; запись $d\mid n$ означает, что $n=dk$ для некоторого целого $k$. НОД и НОК двух натуральных чисел также являются натуральными числами.
\concept{Маршрут обучения} После каждого подраздела выполните ровно три коротких упражнения; изучите схему на отдельной странице; решите десять самостоятельных задач и проверьте ответы. Последняя часть --- обзорная контрольная работа.
\concept{Главное различие} При поиске НОД берём \emph{общую часть}, при поиске НОК --- \emph{все необходимые множители}. Это различие особенно важно в задачах со словами «максимальный размер» и «первое совпадение».

% BLOCK 1
\blocktitle{1}{Что такое НОД и НОК}
\subsection{1.1. Общие делители и общие кратные}
\concept{Цель} Научиться распознавать, какая величина нужна в задаче, ещё до вычислений.
\concept{Определения} \textbf{Наибольший общий делитель} чисел $a,b\in\mathbb N$ --- наибольшее натуральное число, делящее и $a$, и $b$. Обозначение: $\gcdru{a}{b}$. \textbf{Наименьшее общее кратное} --- наименьшее положительное число, кратное обоим числам. Обозначение: $\lcmru{a}{b}$.
\example{Рассмотрим $12$ и $18$. Положительные делители $12$:
\[1,2,3,4,6,12;\]
положительные делители $18$:
\[1,2,3,6,9,18.\]
Наибольший общий --- $6$, поэтому $\gcdru{12}{18}=6$. Положительные кратные $12$: $12,24,36,48,\ldots$; кратные $18$: $18,36,54,\ldots$. Первое совпадение --- $36$, поэтому $\lcmru{12}{18}=36$.}
\concept{Сравнение} Если две верёвки длиной $12$ и $18$ см нужно разрезать на \emph{максимально длинные одинаковые} целые отрезки без остатка, длина равна НОД: $6$ см. Если два сигнала подаются через $12$ и $18$ минут, следующий \emph{общий} сигнал будет через НОК: $36$ минут.
\concept{Свойства из определения} Для натуральных $a,b$:
\[
\gcdru{a}{b}\mid a,\qquad \gcdru{a}{b}\mid b,
\quad a\mid\lcmru{a}{b},\quad b\mid\lcmru{a}{b}.
\]
Также $\gcdru{a}{a}=a$, $\lcmru{a}{a}=a$, $\gcdru{1}{a}=1$ и $\lcmru{1}{a}=a$. Если $a\mid b$, то НОД равен $a$, а НОК --- $b$.
\warning{НОД --- не «наибольшее кратное», а НОК --- не «наименьший делитель». Эти понятия устроены противоположным образом.}
\minitrain
\mini{1}{Найдите $\gcdru{8}{20}$. ($\boldsymbol{4}$.)}
\mini{2}{Найдите $\lcmru{6}{10}$. ($\boldsymbol{30}$.)}
\mini{3}{Найдите $\gcdru{9}{27}$ и $\lcmru{9}{27}$. ($\boldsymbol{9}$ и $\boldsymbol{27}$.)}

\subsection{1.2. Вычисление через простые множители}
\concept{Цель} Научиться находить обе величины по каноническому разложению, не перечисляя делители и кратные.
\concept{Теорема} Пусть $p_1,\ldots,p_k$ --- все простые числа, встречающиеся в разложениях $a$ или $b$ (недостающий показатель принимаем равным нулю):
\[
a=\prod_{i=1}^k p_i^{\alpha_i},\qquad b=\prod_{i=1}^k p_i^{\beta_i}.
\]
Тогда
\[
\boxed{\gcdru{a}{b}=\prod_{i=1}^k p_i^{\min(\alpha_i,\beta_i)}},
\]
\[
\boxed{\lcmru{a}{b}=\prod_{i=1}^k p_i^{\max(\alpha_i,\beta_i)}}.
\]
\concept{Почему так} Общий делитель может содержать простое $p_i$ не больше, чем его содержит каждое из двух чисел. Поэтому берём \emph{меньший показатель}. Общее кратное должно включать достаточно множителей для делимости на оба числа --- выбираем \emph{больший показатель}.
\example{Найдём НОД и НОК чисел $360$ и $840$:
\begin{align*}
360&=2^3\cdot3^2\cdot5,\\
840&=2^3\cdot3\cdot5\cdot7;\\
\gcdru{360}{840}&=2^3\cdot3\cdot5=120,\\
\lcmru{360}{840}&=2^3\cdot3^2\cdot5\cdot7=2520.
\end{align*}
Проверка: $120\mid360$, $120\mid840$, а $2520$ делится на $360$ и на $840$.}
\example{Если $a=2^5\cdot3^2\cdot11$, $b=2^3\cdot5\cdot11^2$, то
\[
\gcdru{a}{b}=2^3\cdot11=88,\qquad
\lcmru{a}{b}=2^5\cdot3^2\cdot5\cdot11^2.
\]
Простое число $3$ отсутствует во втором числе, а $5$ --- в первом, поэтому эти простые множители не входят в НОД.}
\warning{В НОД нельзя переносить простые множители, которые есть лишь в одном числе. При вычислении НОК такие множители, наоборот, обязательны.}
\minitrain
\mini{1}{Найдите $\gcdru{72}{120}$. ($\boldsymbol{24}$.)}
\mini{2}{Найдите $\lcmru{18}{45}$. ($\boldsymbol{90}$.)}
\mini{3}{Найдите НОД и НОК чисел $2^3\cdot5$ и $2\cdot5^2$. ($\boldsymbol{10}$ и $\boldsymbol{200}$.)}

\schemetitle{НОД и НОК по простым множителям}
\begin{center}
\begin{tikzpicture}[node distance=13mm]
\node[st] (s) {Даны $a,b\in\mathbb N$};
\node[pr,below=of s] (a) {Разложите оба числа\\на простые множители};
\node[pr,below=of a] (b) {Выпишите все встречающиеся\\простые числа $p_i$};
\node[de,below=17mm of b] (q) {Ищем НОД?};
\node[pr,text width=43mm,below left=18mm and 19mm of q] (g) {Меньший показатель\\для каждого $p_i$};
\node[pr,text width=43mm,below right=18mm and 19mm of q] (l) {Больший показатель\\для каждого $p_i$};
\node[pr,text width=74mm,below=29mm of q,yshift=-27mm] (m) {Перемножьте полученные степени};
\node[st,below=of m] (e) {НОД либо НОК};
\draw[ar](s)--(a);\draw[ar](a)--(b);\draw[ar](b)--(q);
\draw[ar](q.south west)--node[ann,above left]{да}(g.north east);
\draw[ar](q.south east)--node[ann,above right]{нет}(l.north west);
\draw[ar](g.south)|-(m.west);\draw[ar](l.south)|-(m.east);\draw[ar](m)--(e);
\end{tikzpicture}
\end{center}
\training{1}{От прямого применения определений к работе с большими разложениями.}
\task{1}{Найдите $\gcdru{18}{30}$.}
\task{2}{Найдите $\lcmru{8}{12}$.}
\task{3}{Перечислите все положительные общие делители чисел $24$ и $36$.}
\task{4}{Назовите три наименьших положительных общих кратных чисел $6$ и $9$.}
\task{5}{Найдите НОД и НОК чисел $84$ и $126$.}
\task{6}{Найдите НОД и НОК чисел $180$ и $252$.}
\task{7}{Найдите НОД и НОК чисел $2^5\cdot3^2\cdot5$ и $2^3\cdot3^4\cdot7$. Ответ допустимо дать в виде произведений степеней.}
\task{8}{Найдите НОД и НОК чисел $360$ и $840$.}
\task{9}{Найдите НОД и НОК чисел $2^6\cdot3^3\cdot11$ и $2^4\cdot3^5\cdot5$. Ответ допустимо дать в виде произведений степеней.}
\task{10}{Найдите наименьшее натуральное число, которое делится на $24$, на $35$ и на $9$.}
\answers{1}{\anstable{
1&$6$&6&$36$ и $1260$\\
2&$24$&7&$72$ и $2^5\cdot3^4\cdot5\cdot7$\\
3&$1,2,3,4,6,12$&8&$120$ и $2520$\\
4&$18,36,54$&9&$2^4\cdot3^3$ и $2^6\cdot3^5\cdot5\cdot11$\\
5&$42$ и $252$&10&$2520$
}}
\selfcheck{Порядок ответов в парных заданиях: сначала НОД, затем НОК. В задаче 9 большие показатели не означают, что нужно вычислять все степени десятичным числом.}

% BLOCK 2
\blocktitle{2}{Алгоритм Евклида и линейные комбинации}
\subsection{2.1. Почему алгоритм Евклида работает}
\concept{Цель} Находить НОД больших чисел быстро, даже если их простые множители неизвестны.
\concept{Свойство делимости} Если $d\mid a$ и $d\mid b$, то $d$ делит любое выражение $xa+yb$ с целыми $x,y$. В частности, он делит $a-b$ и $a-qb$. Обратно, общий делитель $b$ и $a-qb$ делит число $(a-qb)+qb=a$. Значит,
\[
\boxed{\gcdru{a}{b}=\gcdru{b}{a-qb}}.
\]
\concept{Теорема (алгоритм Евклида)} Если $a=qb+r$, где $0\le r<b$, то $\gcdru{a}{b}=\gcdru{b}{r}$. Повторяя деление с остатком, мы получаем уменьшающиеся неотрицательные остатки; процесс заканчивается, а \emph{последний ненулевой остаток} равен НОД.
\example{Найдём $\gcdru{5160}{16920}$. Начинаем с большего числа:
\begin{align*}
16920&=3\cdot5160+1440,\\
5160&=3\cdot1440+840,\\
1440&=1\cdot840+600,\\
840&=1\cdot600+240,\\
600&=2\cdot240+120,\\
240&=2\cdot120+0.
\end{align*}
Следовательно, $\gcdru{5160}{16920}=120$. Не пришлось разлагать ни одно из чисел на простые множители.}
\concept{Сравнение} Для $48$ и $18$ уже достаточно двух переходов:
\[48=2\cdot18+12,\qquad18=1\cdot12+6,\qquad12=2\cdot6.\]
Получаем НОД $6$. Для взаимно простых чисел последним ненулевым остатком будет $1$.
\warning{Не путайте последний \emph{ненулевой} остаток с последним остатком вообще: последний всегда равен нулю, если алгоритм завершён.}
\minitrain
\mini{1}{Найдите алгоритмом Евклида $\gcdru{54}{24}$. ($\boldsymbol{6}$.)}
\mini{2}{Найдите $\gcdru{119}{51}$. ($\boldsymbol{17}$.)}
\mini{3}{Найдите $\gcdru{143}{60}$. ($\boldsymbol{1}$.)}

\subsection{2.2. Представление НОД линейной комбинацией}
\concept{Цель} Научиться не только находить НОД, но и записывать его через исходные числа: это помогает доказывать взаимную простоту.
\concept{Теорема Безу} Для любых натуральных $a,b$ существуют целые числа $x,y$ такие, что
\[
\boxed{ax+by=\gcdru{a}{b}.}
\]
Более того, любая целая линейная комбинация $ax+by$ кратна $\gcdru{a}{b}$. Поэтому если для некоторых целых $x,y$ получилось $ax+by=1$, то $a$ и $b$ взаимно просты.
\concept{Идея доказательства} Каждый остаток алгоритма Евклида представляет собой разность двух предыдущих чисел с целым коэффициентом. Подставляя равенства обратно, можно выразить последний ненулевой остаток, то есть НОД, через $a$ и $b$.
\example{Найдём представление НОД чисел $84$ и $30$:
\begin{align*}
84&=2\cdot30+24,\\30&=1\cdot24+6,\\24&=4\cdot6.
\end{align*}
Обратная подстановка:
\[
6=30-24=30-(84-2\cdot30)=3\cdot30-84.
\]
Значит, $\gcdru{84}{30}=6=-1\cdot84+3\cdot30$.}
\example{Докажем, что при любом натуральном $n$ числа $5n+2$ и $3n+1$ взаимно просты. Их общий делитель должен делить
\[
3(5n+2)-5(3n+1)=1.
\]
Единственный положительный делитель единицы --- $1$, поэтому $\gcdru{5n+2}{3n+1}=1$.}
\warning{Теорема Безу разрешает \emph{отрицательные} коэффициенты $x,y$. Требовать, чтобы оба коэффициента были натуральными, нельзя.}
\minitrain
\mini{1}{Представьте $\gcdru{15}{9}$ как $15x+9y$. ($\boldsymbol{3=2\cdot9-15}$.)}
\mini{2}{Могут ли $14x+21y$ при целых $x,y$ равняться $5$? (\textbf{Нет}.)}
\mini{3}{Найдите $\gcdru{4n+1}{7n+2}$ при $n\in\mathbb N$. ($\boldsymbol{1}$.)}

\schemetitle{нахождение НОД алгоритмом Евклида}
\begin{center}
\begin{tikzpicture}[node distance=14mm]
\node[st] (s) {Даны $a\ge b>0$};
\node[pr,below=of s] (c) {Разделите $a$ на $b$\\с остатком: $a=qb+r$};
\node[de,below=19mm of c] (q) {$r=0$?};
\node[pr,text width=40mm,below left=24mm and 15mm of q] (u) {Замените пару\\$a\leftarrow b$, $b\leftarrow r$};
\node[st,text width=35mm,below right=24mm and 15mm of q] (f) {НОД равен $b$};
\draw[ar](s)--(c);\draw[ar](c)--(q);
\draw[ar](q.south west)--node[ann,above left]{нет}(u.north east);
\draw[ar](q.south east)--node[ann,above right]{да}(f.north west);
\draw[ar](u.west)--++(-9mm,0)|-(c.west);
\end{tikzpicture}
\end{center}
\training{2}{Остатки, обратная подстановка и доказательства взаимной простоты.}
\task{1}{Найдите алгоритмом Евклида $\gcdru{56}{42}$.}
\task{2}{Найдите $\gcdru{91}{35}$.}
\task{3}{Найдите $\gcdru{252}{198}$.}
\task{4}{Найдите $\gcdru{347}{85}$.}
\task{5}{Найдите $\gcdru{1235}{865}$.}
\task{6}{Найдите $\gcdru{2024}{748}$.}
\task{7}{Найдите целые $x,y$, удовлетворяющие $18x+30y=\gcdru{18}{30}$. Достаточно одной пары.}
\task{8}{Представьте $\gcdru{221}{153}$ в виде $221x+153y$ с целыми $x,y$.}
\task{9}{Докажите, что для любого $n\in\mathbb N$ числа $5n+2$ и $3n+1$ взаимно просты. В ответе запишите НОД.}
\task{10}{Найдите все возможные значения $\gcdru{7n+5}{5n+3}$ при $n\in\mathbb N$.}
\answers{2}{\anstable{
1&$14$&6&$44$\\
2&$7$&7&$x=2,\ y=-1$\\
3&$18$&8&$17=-2\cdot221+3\cdot153$\\
4&$1$&9&$1$\\
5&$5$&10&$1,2,4$
}}
\selfcheck{В задаче 10 общая разность $5(7n+5)-7(5n+3)=4$, поэтому НОД делит $4$. При чётном $n$ оба числа нечётны; при $n\equiv1\pmod4$ получаем НОД $4$, при $n\equiv3\pmod4$ --- НОД $2$.}

% BLOCK 3
\blocktitle{3}{Взаимно простые числа и связь НОД с НОК}
\subsection{3.1. Взаимная простота и её последствия}
\concept{Определение} Два натуральных числа называются \textbf{взаимно простыми}, если их НОД равен $1$. Например, $8$ и $15$ взаимно просты, хотя $8$ само по себе составное. Взаимная простота относится к \emph{паре чисел}.
\concept{Свойства} Для натуральных $a,b,c$ верны утверждения:
\begin{itemize}
\item соседние натуральные числа взаимно просты: $\gcdru{n}{n+1}=1$;
\item если $\gcdru{a}{b}=1$ и $a\mid bc$, то $a\mid c$ (лемма Евклида);
\item если $\gcdru{a}{b}=1$, то $\gcdru{a}{a+b}=1$;
\item если $\gcdru{a}{b}=1$, то для положительных $m,n$ числа $a^m,b^n$ также взаимно просты.
\end{itemize}
\concept{Обоснование леммы Евклида} Из теоремы Безу имеем $ax+by=1$. Умножим на $c$:
\[acx+bcy=c.\]
Если $a\mid bc$, то обе части суммы слева делятся на $a$, значит, делится и $c$.
\example{Докажем, что $\gcdru{2n+1}{2n+3}=1$. Любой общий делитель двух нечётных чисел делит их разность $2$. При этом общий делитель нечётен, а единственный положительный нечётный делитель числа $2$ --- $1$.}
\example{Пусть $\gcdru{12}{25}=1$ и $12\mid25k$. Тогда по лемме Евклида $12\mid k$. Например, если $25k=300$, то $k=12$.}
\warning{Два составных числа могут быть взаимно простыми. Нельзя заменять условие $\gcdru{a}{b}=1$ словами «оба числа простые».}
\minitrain
\mini{1}{Взаимно ли просты числа $16$ и $27$? (\textbf{Да}.)}
\mini{2}{Найдите $\gcdru{n}{n+1}$ при $n\in\mathbb N$. ($\boldsymbol{1}$.)}
\mini{3}{Если $\gcdru{9}{10}=1$ и $9\mid10k$, что известно о $k$? ($\boldsymbol{9\mid k}$.)}

\subsection{3.2. Произведение НОД и НОК. Обратные задачи}
\concept{Теорема} Для любых натуральных $a,b$ выполняется
\[
\boxed{\gcdru{a}{b}\cdot\lcmru{a}{b}=ab.}
\]
\concept{Доказательство} Для каждого простого $p$ показатель в произведении слева равен $\min(\alpha,\beta)+\max(\alpha,\beta)=\alpha+\beta$. Это показатель $p$ в $ab$. Следовательно, произведения совпадают.
\concept{Следствие} Если $\gcdru{a}{b}=1$, то $\lcmru{a}{b}=ab$. В общем случае НОК можно вычислить через НОД:
\[\lcmru{a}{b}=\frac{ab}{\gcdru{a}{b}}.\]
\example{Известно, что $\gcdru{48}{n}=8$ и $\lcmru{48}{n}=240$. Тогда
\[48n=8\cdot240,\qquad n=\frac{8\cdot240}{48}=40.\]
Проверка обоих исходных условий обязательна: $\gcdru{48}{40}=8$, $\lcmru{48}{40}=240$.}
\concept{Как восстановить пару} Если $d=\gcdru{a}{b}$, то $a=dx$, $b=dy$ и $\gcdru{x}{y}=1$. Если дополнительно $L=\lcmru{a}{b}$, то
\[\boxed{xy=L/d,\qquad \gcdru{x}{y}=1.}\]
То есть нужно разложить $L/d$ на два \emph{взаимно простых} множителя.
\example{Пусть $\gcdru{a}{b}=5$ и $\lcmru{a}{b}=60$. Тогда $a=5x$, $b=5y$, $xy=12$. Подходящие пары $(x,y)$: $(1,12),(3,4),(4,3),(12,1)$. Следовательно, все упорядоченные пары $(a,b)$:
\[(5,60),\ (15,20),\ (20,15),\ (60,5).\]
Пары $(2,6)$ и $(6,2)$ для $x,y$ не годятся, поскольку имеют общий делитель $2$.}
\warning{Равенства $ab=dL$ недостаточно для восстановления пары: найденные числа нужно дополнительно проверить на заданный НОД.}
\minitrain
\mini{1}{Если НОД чисел равен $4$, а НОК --- $60$, чему равно их произведение? ($\boldsymbol{240}$.)}
\mini{2}{Найдите $\lcmru{14}{25}$. ($\boldsymbol{350}$.)}
\mini{3}{Если $\gcdru{18}{n}=6$ и $\lcmru{18}{n}=90$, найдите $n$. ($\boldsymbol{30}$.)}

\schemetitle{восстановление чисел по НОД и НОК}
\begin{center}
\begin{tikzpicture}[node distance=12mm]
\node[st] (s) {Даны $d=\text{НОД}$\\и $L=\text{НОК}$};
\node[de,below=18mm of s] (q) {$d\mid L$?};
\node[pr,text width=42mm,right=24mm of q] (bad) {Таких натуральных\\чисел нет};
\node[pr,below=19mm of q] (a) {Положите $a=dx$, $b=dy$\\и найдите $xy=L/d$};
\node[pr,below=of a] (b) {Рассмотрите пары множителей\\$x,y$ с $\gcdru{x}{y}=1$};
\node[pr,below=of b] (c) {Восстановите $a=dx$, $b=dy$\\и проверьте оба условия};
\node[st,below=of c] (e) {Все подходящие пары};
\draw[ar](s)--(q);
\draw[ar](q.east)--node[ann,above]{нет}(bad.west);
\draw[ar](q.south)--node[ann,right]{да}(a.north);
\draw[ar](a)--(b);\draw[ar](b)--(c);\draw[ar](c)--(e);
\end{tikzpicture}
\end{center}
\training{3}{Взаимная простота, произведение НОД и НОК, подбор числовых пар.}
\task{1}{Найдите $\gcdru{14}{25}$.}
\task{2}{Найдите $\lcmru{14}{25}$.}
\task{3}{Взаимно ли просты числа $35$ и $64$?}
\task{4}{Найдите $\gcdru{2n+1}{2n+3}$ для любого $n\in\mathbb N$.}
\task{5}{Известно: $\gcdru{a}{b}=6$, $\lcmru{a}{b}=180$. Чему равно $ab$?}
\task{6}{Найдите натуральное $n$, если $\gcdru{48}{n}=8$ и $\lcmru{48}{n}=240$.}
\task{7}{Найдите все упорядоченные пары $(a,b)\in\mathbb N^2$, если $\gcdru{a}{b}=5$ и $\lcmru{a}{b}=60$.}
\task{8}{Найдите натуральное $a$, если $\gcdru{18}{a}=6$, $\lcmru{18}{a}=90$.}
\task{9}{Сколько существует \emph{неупорядоченных} пар натуральных чисел $\{a,b\}$ с $\gcdru{a}{b}=12$ и $ab=4320$? Выпишите их.}
\task{10}{Известно: $\gcdru{a}{b}=9$, $\lcmru{a}{b}=540$. Найдите сумму всех различных значений $a+b$.}
\answers{3}{\anstable{
1&$1$&6&$40$\\
2&$350$&7&$(5,60),(15,20),$ $(20,15),(60,5)$\\
3&Да&8&$30$\\
4&$1$&9&$4$: $\{12,360\},\{24,180\},$ $\{36,120\},\{60,72\}$\\
5&$1080$&10&$1080$
}}
\selfcheck{В задаче 10 из $a=9x$, $b=9y$ получаем $xy=60$ и $\gcdru{x}{y}=1$. Четыре неупорядоченных пары: $(1,60),(3,20),(4,15),(5,12)$; их суммы, умноженные на $9$, равны $549,207,171,153$.}

% BLOCK 4
\blocktitle{4}{Прикладные и доказательные задачи}
\subsection{4.1. Группировка, разрезание и периодические события}
\concept{Цель} Научиться превращать сюжет в задачу на делимость, а затем выбирать между НОД и НОК.
\concept{Сигнальные слова} «Наибольшая возможная длина одинаковых отрезков», «максимальное число одинаковых комплектов» обычно приводят к НОД. «Первый общий момент», «минимальный общий период» --- к НОК. Однако окончательный выбор всегда следует проверять по самому условию.
\example{Имеются две ленты длиной $84$ см и $60$ см. Каждую нужно разрезать без остатка на равные между собой отрезки максимальной длины. Пусть длина отрезка $d$ см. Тогда $d\mid84$ и $d\mid60$, следовательно,
\[d\le\gcdru{84}{60}=12.\]
Длина $12$ см действительно подходит: получится $84/12=7$ и $60/12=5$ отрезков, всего $12$.}
\example{Два сигнала звучат через каждые $18$ и $24$ минуты. Сейчас они совпали. Следующее совпадение наступит через число минут, кратное и $18$, и $24$:
\[18=2\cdot3^2,\quad24=2^3\cdot3,\quad\lcmru{18}{24}=2^3\cdot3^2=72.\]
Следовательно, оба сигнала прозвучат вместе через $72$ минуты.}
\concept{Три и более числа} Для нескольких чисел можно аналогично использовать общие простые множители (для НОД --- минимальные показатели, для НОК --- максимальные) или последовательно вычислять НОД/НОК пар чисел.
\warning{Слова «наибольший» и «наименьший» сами по себе не определяют метод. Например, \emph{наибольшее количество наборов} --- НОД, хотя речь идёт о количестве, а не длине.}
\minitrain
\mini{1}{Максимальная длина равных отрезков из лент $45$ см и $60$ см? ($\boldsymbol{15}$ см.)}
\mini{2}{Через сколько минут совпадут сигналы с периодами $6$ и $10$ мин? ($\boldsymbol{30}$ мин.)}
\mini{3}{Наибольшее число одинаковых наборов из $36$ и $48$ предметов двух видов? ($\boldsymbol{12}$ наборов.)}

\subsection{4.2. НОД выражений и нестандартные условия}
\concept{Цель} Научиться ограничивать НОД с помощью суммы или разности выражений и переходить от доказательства к поиску всех значений.
\concept{Главный приём} Любой общий делитель $A(n)$ и $B(n)$ делит все целочисленные линейные комбинации этих выражений. Если удаётся получить небольшую константу $C$, то
\[\gcdru{A(n)}{B(n)}\mid |C|\quad(C\ne0).\]
Это сужает список кандидатов для НОД до положительных делителей числа $|C|$.
\example{Найдём все возможные значения $\gcdru{n+4}{n+10}$. Поскольку разность равна $6$,
\[\gcdru{n+4}{n+10}=\gcdru{n+4}{6}.\]
Значит, возможны только $1,2,3,6$. Все значения достигаются: для $n=3,4,5,2$ соответственно НОД равен $1,2,3,6$.}
\example{Докажем, что $\gcdru{3n+2}{5n+3}=1$ при всех натуральных $n$:
\[5(3n+2)-3(5n+3)=1.\]
Общий делитель обязан делить $1$, поэтому он равен $1$. Такое короткое доказательство сильнее простого перебора значений $n$.}
\example{Для каких натуральных $n$ выполнено $\gcdru{n+4}{n+10}=6$? Достаточно и необходимо, чтобы $6\mid n+4$, поскольку НОД равен $\gcdru{n+4}{6}$. Поэтому
\[n+4=6k,\quad n=6k-4=6(k-1)+2.\]
Среди натуральных $n$ получаем $n=2,8,14,20,\ldots$; это все числа вида $n=6t+2$ при целом $t\ge0$.}
\warning{Из $d\mid6$ нельзя заключать, что $d=6$. Нужно проверить, какие делители константы действительно достигаются.}
\minitrain
\mini{1}{Найдите $\gcdru{n}{n+1}$ при $n\in\mathbb N$. ($\boldsymbol{1}$.)}
\mini{2}{Чему может равняться $\gcdru{n+2}{n+8}$? ($\boldsymbol{1,2,3,6}$.)}
\mini{3}{Найдите наименьшее натуральное $n$, при котором $\gcdru{n+4}{n+10}=6$. ($\boldsymbol{2}$.)}

\schemetitle{выбор метода в текстовой задаче}
\begin{center}
\begin{tikzpicture}[node distance=14mm]
\node[st] (s) {Прочитайте условие};
\node[pr,below=of s] (a) {Запишите условия\\делимости или кратности};
\node[de,below=18mm of a] (q) {Нужно найти НОД?};
\node[pr,text width=43mm,below left=22mm and 17mm of q] (g) {Вычислите НОД};
\node[pr,text width=43mm,below right=22mm and 17mm of q] (l) {Для первого общего кратного вычислите НОК};
\node[pr,text width=78mm,below=65mm of q] (t) {Переведите число обратно\\в единицы задачи};
\node[st,below=of t] (f) {Проверенный ответ};
\draw[ar](s)--(a);\draw[ar](a)--(q);
\draw[ar](q.south west)--node[ann,above left]{да}(g.north east);
\draw[ar](q.south east)--node[ann,above right]{нет}(l.north west);
\draw[ar](g.south)|-(t.west);\draw[ar](l.south)|-(t.east);\draw[ar](t)--(f);
\end{tikzpicture}
\end{center}
\training{4}{Практические ситуации и логика доказательства для выражений с параметром.}
\task{1}{Ленты длиной $84$ и $60$ см режут на одинаковые отрезки без остатка. Найдите наибольшую длину отрезка.}
\task{2}{Два сигнала подаются через каждые $18$ и $24$ минуты. Через сколько минут после одновременного сигнала они снова совпадут?}
\task{3}{Из $72$ карандашей и $96$ блокнотов составляют максимально возможное число одинаковых комплектов без остатка. Сколько комплектов будет и что входит в один?}
\task{4}{Автобусы двух маршрутов отправились вместе в $9{:}00$ и далее отправляются через $20$ и $30$ минут. В какое время они снова отправятся вместе?}
\task{5}{Прямоугольную панель $126\times210$ см выкладывают целыми квадратными плитками максимально возможной стороны, без обрезков. Найдите сторону и количество плиток.}
\task{6}{Из $72$ мальчиков и $90$ девочек формируют максимально возможное число групп с одинаковым составом. Сколько групп и сколько мальчиков и девочек в каждой?}
\task{7}{Найдите наименьшее натуральное $n>20$, при котором $\gcdru{n}{n+14}=7$.}
\task{8}{Докажите, что для всякого $n\in\mathbb N$ числа $3n+2$ и $5n+3$ взаимно просты. В ответе укажите их НОД.}
\task{9}{Сколько натуральных $n\le100$ удовлетворяют равенству $\gcdru{n+4}{n+10}=6$?}
\task{10}{Два робота стоят в начальной клетке кругового пути из $84$ клеток. Каждый минутный шаг первый робот перемещается вперёд на $6$ клеток, второй --- на $10$ клеток. Через сколько минут они впервые одновременно снова окажутся в начальной клетке?}
\answers{4}{\anstable{
1&$12$ см&6&$18$ групп; $4$ мальчика, $5$ девочек\\
2&$72$ мин&7&$21$\\
3&$24$ комплекта; $3$ карандаша и $4$ блокнота&8&$1$\\
4&$10{:}00$&9&$17$\\
5&$42$ см; $15$ плиток&10&$42$ мин
}}
\selfcheck{Для задачи 10 первый робот возвращается в начальную клетку при $84\mid6t$, то есть при $14\mid t$; вторая --- при $84\mid10t$, что эквивалентно $42\mid t$. Первый общий момент: $\lcmru{14}{42}=42$ минуты.}

\clearpage\section*{Итоговая обзорная контрольная работа}
\textcolor{muted}{10 заданий: 3 базовых, 4 среднего уровня, 2 повышенного уровня и 1 комплексное. Решения записывайте обоснованно.}\par\vspace{3mm}\rule{\linewidth}{.4pt}\vspace{2mm}
\task{1}{Вычислите $\gcdru{72}{90}$.}
\task{2}{Вычислите $\lcmru{15}{24}$.}
\task{3}{Найдите НОД чисел $2^4\cdot3^2\cdot7$ и $2^2\cdot3^5\cdot5$.}
\task{4}{Найдите алгоритмом Евклида $\gcdru{414}{198}$.}
\task{5}{Для натуральных $a,b$ известно, что $\gcdru{a}{b}=12$, $\lcmru{a}{b}=420$. Найдите $ab$.}
\task{6}{Найдите наименьшее натуральное $n>50$, для которого $\gcdru{n}{n+9}=9$.}
\task{7}{Найдите наименьшее натуральное число, кратное одновременно $24$, $36$ и $45$.}
\task{8}{Найдите все неупорядоченные пары натуральных чисел с НОД $4$ и НОК $120$.}
\task{9}{Докажите, что для любого $n\in\mathbb N$ числа $4n+1$ и $9n+2$ взаимно просты. В ответе запишите НОД.}
\task{10}{На двух \emph{раздельных} круговых дорожках длиной $168$ и $120$ м размещают отметки с одним и тем же \emph{максимально возможным} шагом, причём отметка есть в начале каждой дорожки. Два спортсмена одновременно стартуют с начальных отметок и бегут с постоянными скоростями $6$ и $5$ м/мин соответственно. Определите: 1) расстояние между отметками; 2) суммарное число отметок на обеих дорожках; 3) первый момент (в минутах после старта), когда оба спортсмена снова окажутся на начальных отметках одновременно.}
\clearpage\section*{Ответы к итоговой обзорной контрольной работе}
\anstable{
1&$18$&6&$54$\\
2&$120$&7&$360$\\
3&$36$&8&$\{4,120\},\{8,60\},$ $\{12,40\},\{20,24\}$\\
4&$18$&9&$1$\\
5&$5040$&10&1) $24$ м; 2) $12$ отметок; 3) $168$ мин
}
\concept{Контроль задачи 8} При $a=4x$, $b=4y$ имеем $xy=30$, $\gcdru{x}{y}=1$. Попарно простые множители: $(1,30),(2,15),(3,10),(5,6)$ без учёта порядка.
\concept{Контроль задачи 9} Равенство $9(4n+1)-4(9n+2)=1$ означает, что общий делитель делит $1$.
\concept{Контроль задачи 10} На дорожках $168/24=7$ и $120/24=5$ отметок. Первый спортсмен возвращается каждые $168/6=28$ мин, второй --- каждые $120/5=24$ мин; $\lcmru{28}{24}=168$ мин.

\clearpage\section*{Памятка: основные факты}
\begin{center}
\small\begin{tabularx}{\linewidth}{>{\raggedright\arraybackslash}p{43mm}X}\toprule
\textbf{Что нужно знать}&\textbf{Формула или смысл}\\\midrule
НОД&Наибольший положительный общий делитель $a$ и $b$\\
НОК&Наименьшее положительное общее кратное $a$ и $b$\\
Разложение на простые&НОД: меньшие показатели; НОК: большие\\
Связь произведений&$\gcdru{a}{b}\cdot\lcmru{a}{b}=ab$\\
Взаимная простота&$\gcdru{a}{b}=1\Longrightarrow\lcmru{a}{b}=ab$\\
Алгоритм Евклида&$\gcdru{a}{b}=\gcdru{b}{r}$, если $a=qb+r$\\
Теорема Безу&$ax+by=\gcdru{a}{b}$ при некоторых целых $x,y$\\
Лемма Евклида&Если $\gcdru{a}{b}=1$ и $a\mid bc$, то $a\mid c$\\
Восстановление пары&$a=dx$, $b=dy$, $xy=L/d$, $\gcdru{x}{y}=1$\\
Прикладные задачи&Максимальная общая длина/число групп --- НОД; первый общий период --- НОК\\\bottomrule
\end{tabularx}
\end{center}
\concept{Самопроверка} Для НОД результат должен делить оба исходных числа. Для НОК оба исходных числа должны делить результат. Для обратной задачи проверьте \emph{оба} условия, а не только равенство $ab=\text{НОД}\cdot\text{НОК}$.
\end{document}